% ===================================================================
%  TABEL Nr. 13 — Hotell. --- Restoran. Kohvik.
%  Traduction 2026 d'apres texto/io, controlee sur texto/fr et
%  texto/fi. Consigne : voir texto/et/10-tabel-01.tex.
%
%  LE POURBOIRE DE L'ALINEA 9 EST LA QUATRIEME REPONSE DE LA SERIE A
%  LA MEME PIECE DE MONNAIE, ET ELLE OBLIGE A PRECISER LA REGLE DU
%  TABLEAU 2.
%
%    tcheque et lituanien : ils CALQUENT le Trinkgeld allemand ;
%    luxembourgeois : il le PREND tel quel ;
%    romanche : il prend l'italien « buonamano » ;
%    ESTONIEN : « jootraha », de « joot » la boisson et « raha »
%      l'argent — UN CALQUE, morpheme pour morpheme.
%
%  OR CETTE LANGUE-CI EST CELLE QUI A REFUSE « velo » AU TABLEAU 2 ET
%  « pompier » AU TABLEAU 11. Elle a fait « jalgratas » la
%  roue-a-pied et « tuletõrjujad » les repousseurs-de-feu, sans rien
%  devoir a personne. Pourquoi cede-t-elle ici ?
%
%  PARCE QU'UNE LANGUE PLANIFIEE REFUSE LA FORME, PAS L'IDEE. Il faut
%  y regarder de pres : « jalgratas » n'est pas un calque de
%  Fahrrad — ce serait « sõiduratas », la roue-de-course — c'est une
%  description independante, faite par quelqu'un qui a regarde
%  l'objet et non le mot allemand. « jootraha », lui, est Trinkgeld
%  trait pour trait.
%
%  LA DIFFERENCE EST DANS L'OBJET. Une bicyclette se decrit de dix
%  facons — par la roue, par les pedales, par la vitesse — et l'on
%  peut choisir. UN POURBOIRE NE SE DECRIT QUE D'UNE : c'est de
%  l'argent qu'on donne pour que l'autre boive. L'allemand n'avait
%  pas invente cette description, il l'avait trouvee ; l'estonien la
%  trouve a son tour et tombe sur la meme.
%
%  ON NE PEUT DONC PAS DIRE D'UN CALQUE QU'IL EST UNE DEPENDANCE
%  TANT QU'ON N'A PAS VERIFIE QU'UNE AUTRE DESCRIPTION ETAIT
%  POSSIBLE. Le tableau 13 luxembourgeois avait fait du calque le
%  signe d'une distance ; l'estonien montre le cas ou il n'est le
%  signe de rien du tout.
% ===================================================================

%%K t13-noto-1 noto
\VUnotes{143.2mm}{%
\textit{\VUgras{(*) Toa üksikasjade kohta vaata tabelit nr.~6.}}%
}

%%K t13-apar-1 apar
\VUsaut{3.0mm}
\VUtitre{15.0pt}{60}{TABEL Nr. 13}
\VUsaut{1.6mm}
\VUpk{10.2pt}{40}{Hotell. --- Restoran.}
\VUsaut{2.0mm}
\VUpk{10.2pt}{40}{Kohvik.}
\VUsaut{3.4mm}

%%K t13-01-1 p
1. --- Olles eile õhtul Royani saabunud, asusin ma elama «\textit{Ookeani hotelli}»,
mida oli mulle soovitanud üks sõber.

%%K t13-01-2 p
Mulle anti ilus \VUgras{tuba} \textsuperscript{(2)} kolmandal \VUgras{korrusel}
\textsuperscript{(1)}, kus ma magasin väga hästi (*).

%%K t13-02-1 p
2. --- Täna hommikul, väljudes oma toast, kuhu \VUgras{toatüdruk} \textsuperscript{(3)}
oli toonud hommikusöögi, astusin ma \VUgras{suitsetamistuppa} \textsuperscript{(5)}, mida
kasutatakse ka \VUgras{lugemistoana} \textsuperscript{(4)}, et lugeda \VUgras{ajalehti}
\textsuperscript{(6)}, suitsetades \VUgras{sigaretti} \textsuperscript{(8)}. Pärast
lugemist läksin ma alla \VUgras{liftiga} \textsuperscript{(9)} ja läksin linnas jalutama.

%%K t13-03-1 p
3. --- Kuna ma olin unustanud küsida hotelli \VUgras{ametnikult} \textsuperscript{(12)},
mis kell antakse sööke \VUgras{ühisel laual} \textsuperscript{(11)}, tulin ma vara tagasi
ja kasutasin seda ära, et minna suplema hotelli \VUgras{vannituppa}
\textsuperscript{(10)}. Siis läksin ma uuesti \VUgras{kassaruumi} \textsuperscript{(15)},
et helistada ühele oma sõbrale. Kuid \VUgras{telefon} \textsuperscript{(13)} oli hõivatud;
nähes mu kimbatust, palus \VUgras{kassapidaja} \textsuperscript{(14)}, kes oli
registreerimas \VUgras{arvet} \textsuperscript{(17)} \VUgras{külaliste raamatusse}
\textsuperscript{(16)}, \VUgras{portjeel} \textsuperscript{(18)} saata \VUgras{pikkolo}
\textsuperscript{(19)} mu asja ajama \VUgras{jalgrattaga} \textsuperscript{(20)}.

%%K t13-04-1 p
4. --- Ma tänasin teda ja läksin välja, et anda jootraha \VUgras{kandjale}
\textsuperscript{(24)}, kes tõi jaamast oma \VUgras{käruga} \textsuperscript{(25)} mu
\VUgras{kübarakarbi} \textsuperscript{(26)}, \VUgras{mu kohvri} \textsuperscript{(27)} ja
\VUgras{mu reisikoti} \textsuperscript{(28)}. \VUgras{Kirjakandja} \textsuperscript{(21)},
kes oli äsja saabunud, andis mulle \VUgras{kirja} \textsuperscript{(22)}.

%%K t13-05-1 p
5. --- Ma jäin veel mõneks ajaks ukse lävele, \VUgras{kirjakasti} \textsuperscript{(23)}
juurde, vaadates liiklust tänaval. \VUgras{Juht} \textsuperscript{(30)}, kes juhib
\VUgras{omnibust} \textsuperscript{(29)}, laadis oma vankrile \VUgras{kohvrit}
\textsuperscript{(31)}, mis kuulub ühele \VUgras{reisijale} \textsuperscript{(32)}, keda
\VUgras{hotelli peremees} \textsuperscript{(33)} saatis. Noor \VUgras{telegrafist}
\textsuperscript{(35)} tõi kiirustamata \VUgras{telegrammi} \textsuperscript{(34)} ühele
hotelli kliendile; \VUgras{turist} \textsuperscript{(36)} oma \VUgras{fotoaparaadiga}
\textsuperscript{(38)} üle õla uuris oma \VUgras{teejuhti} \textsuperscript{(37)}.

%%K t13-tit-1 sub
\VUsaut{4.0mm}
\VUpk{10.2pt}{40}{Lõunasöögi aeg.}
\VUsaut{3.0mm}

%%K t13-06-1 p
6. --- Võre ees tukkus muretu \VUgras{käskjalg} \textsuperscript{(39)} oma
\VUgras{kandekonksu} \textsuperscript{(40)} ja oma \VUgras{kingaviksimise kasti}
\textsuperscript{(41)} vahel, sellal kui noor \VUgras{pesunaine} \textsuperscript{(42)},
kes kandis pesu\VUgras{korvi} \textsuperscript{(43)}, naeris \VUgras{ülemkelneri}
\textsuperscript{(44)} piduliku ilme üle, kes seisis ukse kõrval.

%%K t13-07-1 p
7. --- Nähes \VUgras{kokka} \textsuperscript{(45)}, kes väljus oma \VUgras{köögist}
\textsuperscript{(47)}, mis asub \VUgras{kelderkorrusel} \textsuperscript{(48)}, et saata
\VUgras{poiss} \textsuperscript{(46)} asja ajama, meenus mulle, et lõunasöögi aeg on
lähedal, ja ma astusin restorani, ületades õue, kus \VUgras{autojuht}
\textsuperscript{(50)} pani garaaži oma \VUgras{auto} \textsuperscript{(49)}, sellal kui
\VUgras{mehaanik} \textsuperscript{(52)} parandas, \VUgras{jalgratturi}
\textsuperscript{(53)} juhiste järgi, \VUgras{bensiinikolmerattalist}
\textsuperscript{(51)}, millega oli äsja õnnetus juhtunud.

%%K t13-08-1 p
8. --- Ma küsisin noorelt \VUgras{pikkololt} \textsuperscript{(54)}, kes kandis
\VUgras{õlleklaase} \textsuperscript{(55)}, kas ühine laud on veranda all, ja tema jaatava
vastuse peale võtsin ma koha ühe laua juures ja lasksin endale anda suurepärase söögi.

%%K t13-noto-2 noto
\VUnotes{143.2mm}{%
\textit{\VUgras{(*) Sõnaraamatusse lisatud roogade nimekirja abil saab õpetaja kergesti
muuta alljärgnevat menüüd.}}%
}

%%K t13-tit-2 sub
\VUsaut{4.0mm}
\VUpk{10.2pt}{40}{Suurepärane lõunasöök.}
\VUsaut{3.0mm}

%%K t13-09-1 p
9. --- Kelner tõi mulle lõunasöögi menüü (*) ja veininimekirja. Ma valisin ja tellisin
eelroaks \VUgras{krevetid} \VUgras{võiga} ja, \VUgras{esiroaks}, \VUgras{rupsi Caeni
moodi}. Ma ei söönud \VUgras{kala}, vaid palusin portsjoni lamba\VUgras{kintsu} ja,
\VUgras{juurviljaroaks}, \VUgras{spinatit} koorega. Siis, kuna ükski \VUgras{vaheroog}
mulle ei meeldinud, lasksin ma tuua mitmesuguseid magustoite, \VUgras{juustu},
\VUgras{puuvilju} ja \VUgras{küpsiseid}. Ma lisasin veel suurepärase Saint-Emilioni
\VUgras{veini} ja, oma söögiga väga rahul, lahkusin ma lauast, olles andnud
\VUgras{kelnerile} \textsuperscript{(57)} ilusa \VUgras{jootraha} \textsuperscript{(59)}.

%%K t13-10-1 p
10. --- Nähes mind lahkumiseks valmis, tegi \VUgras{juhataja} \textsuperscript{(60)}
mulle ettepaneku minna rõdule, et imetleda \VUgras{ookeani} \textsuperscript{(62)}
suurepärast panoraami. Kaugel hakkasid paistma kala\VUgras{paadikesed}
\textsuperscript{(63)}, \VUgras{purjelaevad} \textsuperscript{(64)} ja tohutu
\VUgras{pakilaev} \textsuperscript{(65)}, mille must siluett kerkis esile valgetel
\VUgras{pilvedel} \textsuperscript{(67)}, mis on \VUgras{silmapiiril}
\textsuperscript{(66)}. Kerge briis pani lehvima \VUgras{lipu} \textsuperscript{(70)}
\VUgras{sildi} \textsuperscript{(69)} kohal, mis on \VUgras{halli} \textsuperscript{(68)}
küljes.

%%K t13-tit-3 sub
\VUsaut{3.0mm}
\VUpk{10.2pt}{40}{Lõbustused.}
\VUsaut{3.0mm}

%%K t13-11-1 p
11. --- Vaatepilt oli nii huvitav, et ma otsustasin homme kohviku terrassile minna,
võttes kaasa \VUgras{binokli} \textsuperscript{(71)} või \VUgras{pikksilma}
\textsuperscript{(72)}.

%%K t13-12-1 p
12. --- Rõdult lahkudes läksin ma hotelli kohvikusse, et võtta \VUgras{jooki}
\textsuperscript{(73)}, mängides \VUgras{piljardipartiid} \textsuperscript{(75)}.
Peatumata läbisin ma kohvikusaali, kus paljud külastajad mängisid \VUgras{malet}
\textsuperscript{(78)}, \VUgras{kabet} \textsuperscript{(79)}, \VUgras{doominokividega}
\textsuperscript{(80)}, \VUgras{triktrakki} \textsuperscript{(81)} või \VUgras{kaarte}
\textsuperscript{(82)}, ja läksin otse üles \VUgras{piljardisaali}
\textsuperscript{(74)}.

%%K t13-13-1 p
13. --- Kui partii oli lõppenud, läksin ma alla esimesele korrusele ja palusin
kelnerilt \VUgras{ümbrikku} \textsuperscript{(84)}, \VUgras{paberit} \textsuperscript{(83)},
\VUgras{postmarki} \textsuperscript{(85)}, \VUgras{sulepead} \textsuperscript{(87)} ja
\VUgras{tinti} \textsuperscript{(86)}, et kirja kirjutada.

%%K t13-13-2 p
Siis kutsusin ma \VUgras{voorimehe} \textsuperscript{(92)}, kelle \VUgras{tõld}
\textsuperscript{(88)} oli peatunud kohviku ees. Ma küsisin temalt hinda tundide järgi
või ühe sõidu eest, ja kui me hinnas kokku leppisime, avas ta mulle \VUgras{tõllaukse}
\textsuperscript{(89)} ja seadis mu sisse. Ta tõusis tagasi \VUgras{pukki}
\textsuperscript{(91)}, pani hobused \VUgras{piitsaga} \textsuperscript{(93)} kiiresti
liikuma ja me väljusime võluvale jalutuskäigule.

\VUsaut{6.0mm}
\VUfilet{22mm}
